\documentclass[10pt,a4paper]{amsart}
\usepackage[margin=19mm]{geometry}
\usepackage{amsmath,amssymb,mathtools}
\usepackage[scr=rsfs,cal=euler]{mathalfa}
\usepackage{xfrac}
\usepackage[unicode,hidelinks,bookmarks=false]{hyperref}
\newcommand{\ie}{i.e.\ }
\newcommand{\cf}{cf.\ }
\newcommand{\resp}{respectively\ }
\newcommand{\Subsection}{\subsection}
\providecommand{\nobreakdash}{\nobreak\hbox{-}}
\setlength{\emergencystretch}{2em}
\multlinegap0pt
\usepackage{amssymb}
\usepackage{bm}
\usepackage{booktabs}
\usepackage[shortlabels]{enumitem}
\newtheorem{theorem}{Theorem}
\newtheorem{lemma}{Lemma}
\def\Z{\mathbb Z}
\def\eq{\equiv}
\def\f{\frac}
\def\l{\left}
\def\ord{{\rm ord}}
\def\r{\right}
\def\t{\text}
\title{A supercongruence related to Whipple's ${}_5F_4$ formula and Dwork's dash operation}
\author{Chen Wang, He-Xia Ni}
\date{Source: arXiv:2602.13001v1 (CC BY 4.0)}
\begin{document}
\maketitle

\noindent\emph{Source and licence.} A supercongruence related to Whipple's ${}_5F_4$ formula and Dwork's dash operation. Version arXiv:2602.13001v1 (CC BY 4.0), distributed by arXiv under the Creative Commons Attribution 4.0 International licence, http://creativecommons.org/licenses/by/4.0/, as recorded in the arXiv licence metadata for this exact version and shown on the abstract page https://arxiv.org/abs/2602.13001v1. Full licence notice: \texttt{licenses/arxiv-2602.13001-CC-BY-4.0.txt}.\par\smallskip

\noindent\textit{Editorial scope.} Counted unit: the article's lemma evalG (source0.tex lines 543--548). The article's complete original proof of it is delivered with it (source0.tex lines 550--625, one \texttt{proof} environment of 76 lines). No mathematics is written, repaired or re-derived: the statement and the proof are the article's own lines, in the article's own notation, and no retained line is substituted or re-flowed.  Retained uncounted as the source-local context the proof needs: lines 178--184; lines 340--345; lines 363--370; lines 455--458.  Nothing was omitted from any retained span, so the package carries no omission record.  No retained line contains a citation, so the wrapper carries no bibliography and no bibliography entry.  No retained line contains a figure environment, an image-inclusion command or drawing code, so no image is delivered and the package declares no asset dependency.  Editorial formatting only, in addition: an AMS \texttt{amsart} preamble that restates the article's own declarations and macros used by the retained lines, namely the environments \texttt{theorem}, \texttt{lemma} on separate counters, together with the standard packages those lines need.  The retained environments print in this document as Theorem 1, Lemma 1 in order of appearance, whereas the article prints the same items with the numbering of its own full document.  The complete native source of the article as distributed in the arXiv e-print for this exact version is bundled unchanged as \texttt{sources/194603/source0.tex}.
\begin{theorem}\label{mainth1}
Let $c,d,s\in\Z^+$ with $d\geq 2$, $1\leq c,s\leq d$ and $\gcd(cs,d)=1$, and let $p\geq5$ be a prime with $p\eq s\pmod d$. Then, for any $r\in\Z^+$ with $(\f12+\alpha)^{\ast_r}(\f12+\alpha^{\ast_r})\not\eq0\pmod{p}$,  we have
\begin{equation}\label{mainth1eq}
\sum_{k=0}^{p^r-1}(2k+\alpha)\f{(\alpha)_k^3(\f12)_k}{(1)_k^3(\f12+\alpha)_k}\eq \alpha^{\ast_r}p^r -\f{(\alpha^{\ast_r})^3}{(\f12+\alpha)^{\ast_r}}p^{r+2}H_{\alpha^{\ast_r}p-\alpha^{\ast_{r-1}}}^{(2)}\pmod{p^{r+3}},
\end{equation}
where $\alpha=c/d$.
\end{theorem}

\begin{lemma}\label{pochreduce}
For any odd prime $p$, $x\in\Z_p$ and $r\in\Z^+$, we have
$$
(x)_{p^r}=(-1)^rx^{\ast_r}\prod_{j=1}^r\l(p^{p^{j-1}}\f{\Gamma_p(x+p^j)}{\Gamma_p(x)}\r).
$$
\end{lemma}

\begin{lemma}\label{1/2+alpha}
Under the conditions of Theorem \ref{mainth1}, for $l\in\{0,1,2,\ldots,a-1\}$, modulo $p$ we have
\begin{align*}
&\f{(\f12+\alpha)_l}{(\f12+\alpha+p^r)_l}\\
&\qquad\eq \begin{cases}1,\quad&\t{if}\ a<(p^r+1)/2\ \t{or}\ a\geq(p^r+1)/2\ \t{and}\ l<a-(p^r-1)/2,\\
(2{\alpha}^{\ast_r}-1)/(2{\alpha}^{\ast_r}+1),\quad&\t{if}\ a\geq(p^r+1)/2\ \t{and}\ l\geq a-(p^r-1)/2.\end{cases}
\end{align*}
\end{lemma}

\begin{align}
\label{harmonicreduce1eq}p^{2r}\sum_{l=1}^{a}\f{1}{(\alpha+a-l)^2}&\eq p^2H_{\alpha^{\ast_r}p-\alpha^{\ast_{r-1}}}^{(2)}\pmod{p^3},\\
\label{harmonicreduce2eq}p^{2r}\sum_{l=1}^{(p^r-1)/2}\f{1}{(\alpha+a-l)^2}&\eq 0\pmod{p^3}.
\end{align}

\begin{lemma}\label{evalG}
Under the conditions of Theorem \ref{mainth1}, we have
\begin{equation}\label{evalGeq}
\sum_{l=0}^{a-1}G(\alpha+l,p^r)\eq \f{(\alpha^{\ast_r})^3}{(\f12+\alpha)^{\ast_r}}p^{r+2}H_{\alpha^{\ast_r}p-\alpha^{\ast_{r-1}}}^{(2)}\pmod{p^{r+3}}.
\end{equation}
\end{lemma}

\begin{proof}
Clearly,
\begin{align*}
\sum_{l=0}^{a-1}G(\alpha+l,p^r)&=\sum_{l=0}^{a-1}\f{p^{3r}(p^r+2\alpha+2l)}{(\alpha+l)^3}\f{(\alpha+l)_{p^r}^3(\f12)_{p^r}}{(1)_{p^r}^3(\f12+\alpha+l)_{p^r}}\\
&=\f{p^{3r}(\alpha)_{p^r}^3(\f12)_{p^r}}{(1)_{p^r}^3(\f12+\alpha)_{p^r}}\sum_{l=0}^{a-1}\f{(p^r+2\alpha+2l)(\alpha+p^r)_l^3(\f12+\alpha)_l}{(\alpha+l)^3(\alpha)_l^3(\f12+\alpha+p^r)_l}.
\end{align*}
Then, by Lemma \ref{pochreduce}, we have
\begin{align*}
\sum_{l=0}^{a-1}G(\alpha+l,p^r)&=\f{p^{3r}(\alpha^{\ast_r})^3}{2(\f12+\alpha)^{\ast_r}}\prod_{j=1}^r\f{\Gamma_p(\alpha+p^j)^3\Gamma_p(\f12+p^j)\Gamma_p(1)^3\Gamma_p(\f12+\alpha)}{\Gamma_p(\alpha)^3\Gamma_p(\f12)\Gamma_p(1+p^j)^3\Gamma_p(\f12+\alpha+p^j)}\\
&\quad\times \sum_{l=0}^{a-1}\f{(p^r+2\alpha+2l)(\alpha+p^r)_l^3(\f12+\alpha)_l}{(\alpha+l)^3(\alpha)_l^3(\f12+\alpha+p^r)_l},
\end{align*}
where we have used the facts
$$
\l(\f12\r)^{\ast}=\f12\quad \t{and}\quad 1^{\ast}=1.
$$
Since $(\f12+\alpha)^{\ast_r}\not\eq0\pmod{p}$, we have
$$
\ord_p\l(\f{p^{3r}(\alpha^{\ast_r})^3}{(\f12+\alpha)^{\ast_r}}\r)\geq 3r.
$$
It is easy to see that $\ord_p(\alpha+l)\leq r-1$ for $l\in\{0,1,2,\ldots,a-1\}$. It follows that
$$
\ord_p\l(\f{p^{4r}}{(\alpha+l)^3}\r)\geq r+3
$$
and
$$
\ord_p\l(\f{p^{3r}}{(\alpha+l)^2}\r)\geq r+2.
$$
For $k\in\{0,1,2,\ldots,a-1\}$ we have
$$
\f{(\alpha+p^r)_l}{(\alpha)_l}=\prod_{j=0}^{l-1}\f{\alpha+j+p^r}{\alpha+j}=\prod_{j=0}^{l-1}\l(1+\f{p^r}{\alpha+j}\r)\eq1\pmod{p}.
$$
In view of Lemma \ref{1/2+alpha} and the above, we obtain
\begin{equation}\label{evalGkey}
\sum_{l=0}^{a-1}G(\alpha+l,p^r)\eq\f{p^{3r}(\alpha^{\ast_r})^3}{(\f12+\alpha)^{\ast_r}}\sum_{l=0}^{a-1}\f{(\f12+\alpha)_l}{(\alpha+l)^2(\f12+\alpha+p^r)_l}\pmod{p^{r+3}}.
\end{equation}


Below we divide the proof into two cases.

\medskip

\noindent{\it Case 1}. $a<(p^r+1)/2$.

\medskip

By Lemma \ref{1/2+alpha},
$$
\f{(\f12+\alpha)_l}{(\f12+\alpha+p^r)_l}\eq1\pmod{p}.
$$
This, together with \eqref{evalGkey}, gives
\begin{align*}
\sum_{l=0}^{a-1}G(\alpha+l,p^r)&\eq \f{p^{3r}(\alpha^{\ast_r})^3}{(\f12+\alpha)^{\ast_r}}\sum_{l=0}^{a-1}\f{1}{(\alpha+l)^2}\\
&=\f{p^{3r}(\alpha^{\ast_r})^3}{(\f12+\alpha)^{\ast_r}}\sum_{l=1}^{a}\f{1}{(\alpha+a-l)^2}\pmod{p^{r+3}}.
\end{align*}
Then, by \eqref{harmonicreduce1eq}, we obtain \eqref{evalGeq} in this case.

\medskip

\noindent{\it Case 2}. $a\geq(p^r+1)/2$. 

\medskip

By Lemma \ref{1/2+alpha},
$$
\f{(\f12+\alpha)_l}{(\f12+\alpha+p^r)_l}\eq\begin{cases}1\pmod{p},\quad&\t{if}\ l<a-(p^r-1)/2,\\ (2\alpha^{\ast_r}-1)/(2\alpha^{\ast_r}+1)\pmod{p},\quad&\t{if}\ l\geq a-(p^r-1)/2.\end{cases}
$$
Substituting this into \eqref{evalGkey} and using \eqref{harmonicreduce2eq} we have
\begin{align*}
\sum_{l=0}^{a-1}G(\alpha+l,p^r)&\eq \f{p^{3r}(\alpha^{\ast_r})^3}{(\f12+\alpha)^{\ast_r}}\sum_{l=0}^{a-(p^r+1)/2}\f{1}{(\alpha+l)^2}+\f{p^{3r}(\alpha^{\ast_r})^3(2\alpha^{\ast_r}-1)}{(\f12+\alpha)^{\ast_r}(2\alpha^{\ast_r}+1)}\sum_{l=a-(p^r-1)/2}^{a-1}\f{1}{(\alpha+l)^2}\\
&=\f{p^{3r}(\alpha^{\ast_r})^3}{(\f12+\alpha)^{\ast_r}}\sum_{l=(p^r+1)/2}^{a}\f{1}{(\alpha+a-l)^2}+\f{p^{3r}(\alpha^{\ast_r})^3(2\alpha^{\ast_r}-1)}{(\f12+\alpha)^{\ast_r}(2\alpha^{\ast_r}+1)}\sum_{l=1}^{(p^r-1)/2}\f{1}{(\alpha+a-l)^2}\\
&\eq \f{p^{3r}(\alpha^{\ast_r})^3}{(\f12+\alpha)^{\ast_r}}\sum_{l=1}^{a}\f{1}{(\alpha+a-l)^2}\pmod{p^{r+3}}.
\end{align*}
With the help of \eqref{harmonicreduce1eq}, we obtain the desired result.

The proof of Lemma \ref{evalG} is now complete.
\end{proof}

\end{document}
